\section{Symmetries and degenerate indices}
\label{sec:degenerate}

The cubic equations have a useful symmetry that becomes visible after
replacing four indices by two triples. Associate to $(g,h,k,l)$ the pair
\begin{equation}\label{tail:triples}
 \boldsymbol a=(0,k-g,l-h),\qquad
 \boldsymbol c=(g+h,-k,-l).
\end{equation}
The pair is \emph{balanced}, meaning that its six entries sum to zero.
Its cross sums are
\begin{equation}\label{tail:crossmatrix}
 (a_i+c_j)_{i,j=0}^2=
 \begin{pmatrix}
 g+h&-k&-l\\
 h+k&-g&k-g-l\\
 g+l&l-h-k&-h
 \end{pmatrix}.
\end{equation}
We call the pair \emph{nondegenerate} when all nine cross sums are nonzero.
Every balanced pair has the form \eqref{tail:triples} after an opposite
translation $(a_i,c_j)\mapsto(a_i-a_0,c_j+a_0)$: take
$k=-c_1$, $l=-c_2$, $g=-c_1-a_1$, and $h=-c_2-a_2$ after translating.

For a balanced pair, set
\[
 S(\boldsymbol a,\boldsymbol c)
 =\sum_m\prod_{i=0}^2W_{m+a_i}\prod_{j=0}^2W_{c_j-m},
 \qquad
 P(\boldsymbol a,\boldsymbol c)=\prod_{i,j=0}^2W_{a_i+c_j}.
\]

\begin{lemma}[Factorization and permutation symmetry]\label{tail:factor-lemma}
For a nondegenerate pair,
\begin{equation}\label{tail:factor}
 \mathcal C_{g,h}(k,l)
 =\frac{W_{-g-h}W_gW_h}{\delta^5}
       \bigl(\delta^2S(\boldsymbol a,\boldsymbol c)
                      -\tau P(\boldsymbol a,\boldsymbol c)\bigr).
\end{equation}
Its vanishing is preserved by independent permutations of the two
triples and by opposite translation.
\end{lemma}

\begin{proof}
Substitute \eqref{tail:uniform-core} in \eqref{tail:C} and use
$W_aW_{-a}=\delta$ for $a\ne0$. Nondegeneracy excludes every origin
correction. The prefactor in \eqref{tail:factor} is nonzero.
Both $S$ and $P$ are invariant under the permutations; opposite
translation is absorbed by reindexing $m$ in $S$.
\end{proof}

\begin{lemma}[Degenerate cubics]\label{tail:collisions}
If any cross sum in \eqref{tail:crossmatrix} is zero, then
$C(g,h,k,l)=0$ follows from $Q=0$, the coefficient formula
\eqref{tail:uniform-core}, $W_0=-1$, the reciprocal product identities,
and $\delta^2=n\delta+1$.
\end{lemma}

\begin{proof}
We first treat $g,h,g+h\ne0$, then the three boundary cases.

\emph{A zero cross sum with $g,h,g+h\ne0$.}
Keeping the origin corrections in the preceding substitution gives
\begin{equation}\label{tail:balanced}
 \mathcal C_{g,h}(k,l)
 =\frac{W_{-g-h}W_gW_h}{\delta^5}
 \left(\delta^2S-\tau P-\tau^3\delta^3\,
 \ind{\{c_0,c_1,c_2\}=\{-a_0,-a_1,-a_2\}}\right).
\end{equation}
Here equality in the indicator means equality of multisets.
The only possible origin corrections occur at
$(k,l)=(0,-g)$ and $(-h,0)$. These are exactly the two possible
multiset matches: a matching permutation must avoid the three
nonzero diagonal cross sums and hence is one of the two derangements.
A repeated row would also allow a forbidden diagonal match, so a full
match has three distinct rows. In that case,
\[
 S=\delta^3(n-3)+3\delta^2,\qquad P=-\delta^3.
\]
Thus the bracket in \eqref{tail:balanced} vanishes by
$n=\delta-\delta^{-1}$ and $\tau=\delta-1$.

If there is no full match, permute the triples so that $a_0+c_0=0$.
Then
\[
 W_{m+a_0}W_{-m-a_0}=\delta-\tau\ind{m=-a_0}.
\]
Put $H=a_1+c_1$, $V=a_2+c_1$. Both are nonzero: either zero
would force the remaining entries to match as well. The unrestricted
sum of the other four factors is
\[
\sum_tW_tW_{H-t}W_{-t-V}W_{t+V-H}
       =\delta\tau^2Q(-V,H)=0.
\]
Indeed, with $m=t+V$, the sum is
$\delta\sum_m\mathsf B_{m-V,H}\mathsf B_{H,m}$.
The uniform coefficient formula has no origin correction in these
entries because $H\ne0$. The substitution uses
$W_HW_{-H}=\delta$, and $H,V\ne0$ remove both indicator terms
in $Q(-V,H)$.
Writing
$T=\prod_{i=1}^2W_{a_i-a_0}\prod_{j=1}^2W_{c_j+a_0}$,
we obtain $S=-\tau T$. The four interior factors of $P$ pair
to $\delta^2$, and $W_0=-1$, so $P=-\delta^2T$.
Again $\delta^2S-\tau P=0$.

\emph{The boundary cases.}
Set $c=\tau^{-1}$. For $h\ne0$, the coefficient formula gives
\begin{equation}\label{tail:two-products}
 \mathsf A_{u,h}\mathsf A_{h,u}=
 \begin{cases}
 c^2,&u=0\text{ or }u=h,\\
 \delta c^2,&u\ne0,h,
 \end{cases}
 \qquad \delta c^2-c=c^2.
\end{equation}
For $g=0$, $h\ne0$, use the fixed row
$\mathsf A_{0,a}=\ind{a=0}-c$ and $Q=0$ to reduce the cubic sum to
\[
 \mathsf A_{-k,h}\mathsf A_{h,l-k}-c\ind{l=0}.
\]
For $l\ne0$, the uniform formula has no origin corrections and gives
\[
 \mathsf B_{-k,h}\mathsf B_{h,l-k}
 =\delta^{-1}W_{-k}W_{h+k}W_{l-k-h}W_{k-l}
 =\mathsf B_{l,k}\mathsf B_{h+k,l}.
\]
Dividing by $\tau^2$ gives the required equality.
For $l=0$, the target product is
$(\ind{k=0}-c)(\ind{k=-h}-c)$.
Both sides equal $c^2-c$ at $k=0,-h$, and $c^2$ elsewhere,
by \eqref{tail:two-products}.
The case $h=0$, $g\ne0$ follows by interchanging $(g,k)$ and $(h,l)$.
For $h=-g\ne0$, the same quadratic reduction gives
\[
 \mathsf A_{g,k}\mathsf A_{-l,g}-c\ind{k=g+l}.
\]
Substitution treats $k\ne g+l$, and \eqref{tail:two-products}
treats $k=g+l$.

Finally, when $g=h=0$, the cubic sum is
\[
 \ind{k=l=0}
 -c\bigl(\ind{k=0}+\ind{l=0}+\ind{k=l}\bigr)
 +3c^2-nc^3.
\]
The product $\mathsf A_{l,k}\mathsf A_{k,l}$ equals
$(1-c)^2$, $c^2$, or $\delta c^2$, according as all, exactly one,
or none of the equalities $k=0$, $l=0$, $k=l$ hold.
Using $\delta^2=n\delta+1$ in each case proves that the displayed
sum is $\mathsf A_{l,k}\mathsf A_{k,l}-\delta^{-1}$, as required.
\end{proof}

\section{Polynomial interpolation and index reduction}
\label{tail:completion}

The contour calculation gives only a range of Fourier coefficients.
The permutation symmetry now supplies enough zeros to recover the rest.

\subsection{Induction in the ordered range}

\begin{proposition}\label{tail:ordered}
For every odd $n\ge3$, the cubics vanish in the following ordered range:
\[
 \mathcal C_{g,h}(k,l)=0\quad
 (g,h\ge1,\ g+h<n,\ 1\le l\le h,\ k\in G).
\]
\end{proposition}

\begin{proof}
Fix $n$ and use strong induction on the positive integer $g$. By
Proposition~\ref{an:restricted}, for each ordered $g,h,l$, precisely
$n-g+1$ of the $n$ Fourier coefficients of the sequence
$k\mapsto\mathcal C_{g,h}(k,l)$ are known to vanish.
The omitted frequencies form a consecutive cyclic block of
length $g-1$: the principal integers $v=n\theta_R/\pi$,
of parity $g+1$, advance by two as the Fourier index advances
by one, and the frequency condition is $|v|\le n-g$.
For $g=1$ there are no omitted frequencies, proving the base case.
For $g>1$, Fourier inversion therefore gives, with
$\zeta=\exp(2\pi i/n)$,
\begin{equation}\label{tail:polynomial}
 \mathcal C_{g,h}(k,l)=\zeta^{ak}f(\zeta^k),\qquad
                         \deg f\le g-2
\end{equation}
for an integer $a$ and a complex polynomial $f$.
More explicitly, if the omitted block starts at $a$, Fourier inversion
gives
\[
 f(X)=\frac1n\sum_{j=0}^{g-2}
          \widehat{\mathcal C}_{g,h}(a+j,l)X^j,
\]
where the hat denotes the Fourier transform in the first argument $k$.

We show that $k=1,\ldots,g-1$ are zeros.  A degenerate index choice is
already covered by Lemma~\ref{tail:collisions}.  Otherwise
perform the following changes of cubic labels:
\begin{equation}\label{tail:descent}
 (g,h,k,l)\longmapsto
 \begin{cases}
 (k,h,g,l+g-k),&l+g-k\le h,\\
 (k,n-h-k,k+h-l,n-g-l),&l+g-k>h.
 \end{cases}
\end{equation}
These changes arise from permutations and opposite translation
of \eqref{tail:triples}.  With positions numbered $0,1,2$,
the first uses rows $(1,0,2)$ and unchanged columns; the second
uses rows $(2,0,1)$ and columns $(2,1,0)$.  In both cases
translate the first row entry to zero.  Substitution gives
exactly \eqref{tail:descent}, with all entries interpreted
modulo $n$ where necessary.

The new first index is $k<g$.  Both new quadruples are ordered.
This is immediate in the first case; in the second,
\[
 k+(n-h-k)=n-h<n,\qquad
 1\le n-g-l\le n-h-k,
\]
where the last inequality follows from $l+g-k>h$.
The induction hypothesis and the nonzero factor in
\eqref{tail:factor} therefore give each required zero of the
original cubic.  The $g-1$ points
$\zeta,\ldots,\zeta^{g-1}$ are distinct, including for
composite $n$, so \eqref{tail:polynomial} forces $f=0$.
\end{proof}

\subsection{Putting the indices in the ordered range}

\begin{lemma}\label{tail:normalform}
Let $X_{ij}\in\{1,\ldots,n-1\}$ represent the cross sums of a
nondegenerate balanced pair of triples. If
$X_{00}+X_{11}+X_{22}=2n$, independent permutations and opposite
translation put the pair in the ordered range of
Proposition~\ref{tail:ordered}.
\end{lemma}

\begin{proof}
There exist $i\ne j$ with $X_{ij}\ge X_{jj}$.  Otherwise
each diagonal entry would be a strict maximum in its column.
The three row labels would then be distinct.  Let $d_j>0$
be the forward cyclic gap from $a_j$ to the next row label.
Since the other entry in column $j$ is smaller and no cross
sum is zero, $X_{jj}+d_j>n$.  The gaps sum to $n$, so summing
would give $\sum_jX_{jj}>2n$, a contradiction.

Choose such $i,j$, let $r$ be the remaining position, and use
row and column order $(i,r,j)$.  Normalize the first row to
zero.  The resulting cubic indices have
\[
 g=n-X_{rr},\qquad h=n-X_{jj},\qquad l=n-X_{ij}.
\]
Thus $g,h\ge1$, $1\le l\le h$, and
$g+h=2n-X_{rr}-X_{jj}=X_{ii}<n$, as required.
\end{proof}

\section{Reflection by matrix inversion}
\label{sec:reflection}

Matrix inversion transfers a complete cubic block with first indices
$(g,h)$ to the block with first indices $(-h,-g)$. This is the final
step needed to cover every index choice.

\begin{lemma}\label{allodd:reflection}
Suppose the matrix \eqref{foundation:core} satisfies the quadratic equations
and $W_aW_{-a}=\delta$ for $a\ne0$.  Fix $g,h\in G$, put $t=g+h$, and assume $g,h,t\ne0$.  If $C(g,h,k,l)=0$ for every $k,l$, then
$C(-h,-g,k,l)=0$ for every $k,l$.
\end{lemma}

\begin{proof}
We invert the two sides of the cubic block separately. All matrices
have rows and columns indexed by $G$.

\emph{1. Column inverses.} Write $u_p(a)=\mathsf B_{a,p}$ for
$p\ne0$.  The coefficient formula and the quadratic equations give
\begin{equation}\label{allodd:column-inverses}
 u_p(a)=\frac{W_aW_{p-a}}{W_p}=u_p(p-a),\qquad
 \sum_a u_p(a)u_{-p}(x-a)=\tau^2\ind{x=0}.
\end{equation}
For the second identity, multiply $Q(x,p)=0$ by $\tau^2$,
use $\mathsf B_{p,a}=u_{-p}(a-p)=u_{-p}(-a)$, and reindex.
Consequently the circulant with entries $u_p(k-l)$ has inverse
$\tau^{-2}u_{-p}(k-l)$.

\emph{2. The left side.} Define
\[
 L_{g,h}(k,l)=\sum_m u_t(m)\mathsf B_{g,m+k}
                                      \mathsf B_{h,m+l},
 \qquad
 R_{g,h}(k,l)=\mathsf B_{g+l,k}\mathsf B_{h+k,l}.
\]
The assumed cubics say $L_{g,h}=\tau R_{g,h}$.
Let $H_p(k,m)=\mathsf B_{p,k+m}$ and let $D_t$ be the
diagonal matrix with entries $u_t(m)$.  Equation
\eqref{allodd:column-inverses} gives
\[
 H_p^{-1}(k,m)=\tau^{-2}u_p(k+m),\qquad
 L_{g,h}=H_gD_tH_h.
\]
All entries of $D_t$ are nonzero.  Thus $L_{g,h}$ is invertible.
The reciprocal product identities, including the two exceptional samples $m=0,t$,
give
\begin{equation}\label{allodd:diagonal-inverse}
 \frac1{u_t(m)}=\delta^{-1}u_{-t}(-m)
       -\frac{\tau}{\delta}
                    (\ind{m=0}+\ind{m=t}).
\end{equation}
Set
\[
 E(k,l)=u_h(k)u_g(l)+u_h(k+t)u_g(l+t),
 \qquad
 L^{\mathrm r}(k,l)=L_{-h,-g}(-k,-l).
\]
On inserting \eqref{allodd:diagonal-inverse} into
$L_{g,h}^{-1}=H_h^{-1}D_t^{-1}H_g^{-1}$ and replacing
$m$ by $-m$, we obtain
\begin{equation}\label{allodd:left-inverse}
 L_{g,h}^{-1}=\delta^{-1}\tau^{-4}
                         (L^{\mathrm r}-\tau E).
\end{equation}
Here $u_h(k-m)=\mathsf B_{-h,m-k}$, by
\eqref{allodd:column-inverses} and order-three symmetry.

\emph{3. The right side.} We prove, using only the quadratics and
reciprocal products, the corresponding identity
\begin{equation}\label{allodd:right-inverse}
 R^{\mathrm r}=\delta\tau^2 R_{g,h}^{-1}+E,
 \qquad
 R^{\mathrm r}(k,l)=R_{-h,-g}(-k,-l).
\end{equation}
The details below retain the origin corrections in the coefficient formula.
Put
\[
 d_1(k)=W_{h+k}W_{-k},\qquad d_2(l)=W_{g+l}W_{-l},
 \qquad b=u_t(g)=u_t(h)=\frac{W_gW_h}{W_t}\ne0.
\]
Direct substitution of
$\mathsf B_{a,b}=\delta^{-1}W_aW_{b-a}W_{-b}
+\tau^2\delta^{-1}\ind{a=b=0}$ yields
\begin{equation}\label{allodd:right-factor}
 R_{g,h}=\delta^{-2}W_{-t}\,
                   \operatorname{diag}(d_1)M\operatorname{diag}(d_2),
\end{equation}
where
\[
 M(k,l)=u_{-t}(k-l-g)
       +\tau^2(\ind{k=0,l=-g}+\ind{k=-h,l=0}).
\]
The uncorrected matrix $K(k,l)=u_{-t}(k-l-g)$ has inverse
$K^{-1}(k,l)=\tau^{-2}u_t(k-l+g)$.
Writing $U=(e_0,e_{-h})$ and $V=(e_{-g},e_0)$, we have
$M=K+\tau^2 UV^{\mathsf T}$ and
\[
 I_2+\tau^2V^{\mathsf T}K^{-1}U=
       \begin{pmatrix}0&b\\b&0\end{pmatrix}.
\]
This matrix is invertible.  The rank-two inverse formula therefore
proves that $M$, hence $R_{g,h}$, is invertible and that
\begin{equation}\label{allodd:woodbury}
 M^{-1}(k,l)=\tau^{-2}(U_0-U_1-U_2),
\end{equation}
where, for the entry $(k,l)$ in question,
\[
 U_0=u_t(k-l+g),\qquad
 U_1=b^{-1}u_t(k+g)u_t(g-l),\qquad
 U_2=b^{-1}u_t(k+t)u_t(-l).
\]

To check the exceptional entries in
\eqref{allodd:right-inverse}, define
\[
 q_a=W_aW_{-a}=\delta-\tau\ind{a=0},\quad
 \alpha=\frac{q_kq_l}{\delta^2},\quad
 \beta=\frac{q_{g+k}q_{h+l}}{\delta^2},\quad
 Z=\frac{W_{-t}d_2(k)d_1(l)}{\delta^3}.
\]
To see the endpoint corrections explicitly, expand
$R^{\mathrm r}(k,l)=\mathsf B_{-h-l,-k}\mathsf B_{-g-k,-l}$
using \eqref{tail:uniform-core}. The reciprocal products give
\[
\begin{aligned}
 ZR^{\mathrm r}(k,l)
   &=\alpha\beta U_0+\frac{\tau^2}{\delta^2}
       (\ind{k=0,l=-h}+\ind{k=-g,l=0}),\\
 Z\,u_h(k)u_g(l)&=\alpha U_1,\\
 Z\,u_h(k+t)u_g(l+t)&=\beta U_2.
\end{aligned}
\]
In the first line the two corrections come from the two possible
origin entries of $\mathsf B$; they cannot occur together since
$g,h\ne0$. The other two lines follow from $W_tW_{-t}=\delta$
and the definitions of $U_1,U_2$.
Subtracting the two terms of $E$ gives the exact entry identity
\begin{equation}\label{allodd:endpoints}
 Z(R^{\mathrm r}-E)(k,l)
  =\alpha\beta U_0-\alpha U_1-\beta U_2
     +\frac{\tau^2}{\delta^2}
        (\ind{k=0,l=-h}+\ind{k=-g,l=0}).
\end{equation}
If $\alpha\ne1$, then $k=0$ or $l=0$, and $U_1=U_0$.
If $\beta\ne1$, then $k=-g$ or $l=-h$, and $U_2=U_0$.
Both exceptional conditions can hold only at $(k,l)=(0,-h)$
or $(-g,0)$: the assumptions $g,h\ne0$ exclude the other
intersections.  At either point,
$\alpha=\beta=\delta^{-1}$ and $U_0=U_1=U_2=-1$.
Thus the last term of \eqref{allodd:endpoints} cancels
$(\alpha-1)(\beta-1)U_0=-\tau^2/\delta^2$.
In every case the right side of \eqref{allodd:endpoints} is
$U_0-U_1-U_2$.  Equations \eqref{allodd:right-factor} and
\eqref{allodd:woodbury} now give
\[
 Z(R^{\mathrm r}-E)(k,l)
   =\tau^2M^{-1}(k,l)
   =Z\delta\tau^2R_{g,h}^{-1}(k,l).
\]
The factor $Z$ is nonzero, proving
\eqref{allodd:right-inverse} at all entries.

\emph{4. Comparing the inverses.} Now $L_{g,h}=\tau R_{g,h}$ and
\eqref{allodd:left-inverse} imply
\[
 L^{\mathrm r}=\delta\tau^3R_{g,h}^{-1}+\tau E
             =\tau R^{\mathrm r},
\]
where the second equality is \eqref{allodd:right-inverse}.
These are precisely all the cubics with first indices $(-h,-g)$.
\end{proof}

\begin{proposition}\label{allodd:all-cubics}
For every odd $n\ge3$, the analytic coefficient matrix satisfies
$C(g,h,k,l)=0$ for all $g,h,k,l$.
\end{proposition}

\begin{proof}
The degenerate indices, including the cases $g=0$, $h=0$, or $g+h=0$ have
already been proved in Lemma~\ref{tail:collisions}.
Take positive representatives $1\le g,h<n$ with $g+h<n$.
For every nondegenerate pair associated with $(g,h,k,l)$, the
distinguished cross diagonal has representatives
\[
 (g+h,n-g,n-h),
\]
whose sum is $2n$.  The proof of Lemma~\ref{tail:normalform}
therefore puts this pair into the ordered range.  Proposition~\ref{tail:ordered} and
the permutation and translation invariance in
Lemma~\ref{tail:factor-lemma} prove all these cubics, for every $k,l$.

If instead $g+h>n$, the positive representatives
$g'=n-h$, $h'=n-g$ satisfy $g'+h'<n$.  All cubics for $(g',h')$
have just been proved; Lemma~\ref{allodd:reflection} then gives
all cubics for $(-h',-g')=(g,h)$.  This exhausts the indices.
\end{proof}

